\section{The NbS typology library}
\label{sec:typologies}

\subsection{Two levels: cooling archetypes and catalogue entries}
\label{sec:archetypes}

From version \methver{} the library has \textbf{two levels}, and the separation
between them is the central design decision of this release.

A \textbf{cooling archetype} is a cited evidence class. It carries every
performance value --- base cooling score, temperature reduction envelope,
evidence confidence rating, suitability conditions, default co-benefit levels
--- together with the citations behind them. There are \num{18}.

A \textbf{typology} is one of the \num{110} curated entries of the UNEP
\emph{Nature for Cooling} catalogue. It carries identity, family, availability,
the primary cooling mechanism the catalogue ascribes to it, and the one
archetype it inherits. It carries \textbf{no performance value of its own}.

\begin{keynote}
\noindent\textbf{The reason is worth stating without euphemism: solution-specific
cooling literature does not exist for \num{110} typologies.} Nobody has
separately measured a felt-pocket living wall, a suspended-pavement tree pit,
and a green parking lane.

The alternative to archetypes was inventing \num{110} envelopes. That would have
manufactured precision the source catalogue does not contain, and it would have
broken the evidence rule the tool exists to honour --- that no performance value
ships without a citation (\Cref{sec:evidence-rules}).

Under the archetype model an entry inherits a \emph{cited} envelope, and every
result \textbf{names the evidence class it inherited from}. A Miyawaki forest
reports its cooling on the dense-canopy evidence, and says so, rather than
implying a measurement of Miyawaki forests that does not exist.
\end{keynote}

The catalogue's \num{243} entries were curated to \num{110} --- \num{88} merged,
\num{45} dropped --- on the source document's own classification notes. Nearly
every drop is an entry the document itself disqualifies as not an intervention:
a land-use context, a management objective, an umbrella category, a planning
model. Nothing is lost by dropping them. A land-use context is already carried
by the \fieldname{land_use} input, and a strategy is a \emph{package} of kept
entries, which \Cref{sec:packages} makes expressible for the first time.

The \num{110} entries fall into fourteen catalogue families --- tree-based,
street, park, public space, woodland and forest, ground level, water landscape,
ecological network, district, green roof, vertical greening, building,
vegetated shelter, and productive landscape --- and are classified by the source
document as \num{53} \emph{elements} (a discrete thing that is built) and
\num{57} \emph{composites} (a coherent spatial system made of elements). Neither
classification enters a formula; both are reported, because a reader comparing
the tool's menu with the catalogue needs to see which is which.

\subsection{The metric, stated once and enforced throughout}
\label{sec:metric}

\begin{keynote}
\noindent All adopted temperature values in this library are
\textbf{daytime, pedestrian-level \emph{air} temperature reductions}, relative
to an equivalent unimproved site.

They are \textbf{not} land surface temperatures, and they are \textbf{not}
thermal comfort indices such as the physiologically equivalent temperature
(PET) or the universal thermal climate index (UTCI). These are different
physical quantities, measured differently, with different magnitudes and
different policy meanings.

Where a source reports surface temperature or a comfort index, that is stated
explicitly, and the value is used only for the \emph{relative ranking} of
intervention families or to support a caveat. It is never converted into a
\si{\degreeCelsius} air-temperature claim.
\end{keynote}

This discipline is not pedantry. Conflating the three metrics is the single
most common way screening instruments mislead their users, and the divergence
is large: \textcite{venter2021} report an approximately sixfold difference
between surface and canopy-layer heat island magnitudes, and
\textcite{zolch2016} report comfort improvements from trees
(\SI{13}{\percent} PET reduction) that have no direct air-temperature
equivalent. A decision-maker reading ``\qty{3}{\degreeCelsius} of cooling''
assumes air temperature. The methodology gives them air temperature, and where
it cannot, it says so.

Two of the archetypes retrieved for this version exist \emph{because} a source
measured both quantities on the same plots and the two diverged by an order of
magnitude (\Cref{sec:non-canopy,sec:shade-structure}). The discipline is
therefore not only a presentational rule; it changed the library.

A direct consequence is that the tool systematically \emph{understates} the
benefit of shade-dominant interventions, because radiant load reduction --- the
mechanism through which shade most improves human thermal comfort --- does not
appear in an air-temperature figure. This is recorded as a limitation in
\Cref{sec:limit-airtemp} rather than corrected by an unsourced adjustment.

\subsection{The eighteen archetypes}
\label{sec:adopted}

\Cref{tab:typologies} gives the adopted values for all eighteen cooling
archetypes and the number of catalogue entries that inherit each. Derivations
are summarised in \Cref{sec:derivation} and reproduced in full, with citations,
in \Cref{app:typologies}.

\begin{table}[htbp]
\centering
\caption[Adopted archetype values]{The eighteen cooling archetypes with their
adopted values at methodology version \methver{}. All temperature envelopes are
daytime, pedestrian-level \emph{air} temperature reductions in degrees Celsius.
The base cooling score is a \numrange{0}{100} index expressing relative
expected cooling performance before site adjustment; it is not a temperature.
Evidence confidence is the rating defined in \Cref{sec:confidence-levels}.
\emph{Provenance} records where the archetype's values come from: \emph{v1.1}
means the value is the previously published typology value, unchanged;
\emph{new} means the evidence was retrieved for this version; \emph{derived}
means no class-specific evidence exists and the envelope is bounded from an
adjacent class with the bounding argument stated. \emph{Entries} is the number
of the \num{110} catalogue entries inheriting that archetype.}
\label{tab:typologies}
\small
\begin{tabular}{@{}llrclr@{}}
\toprule
\textbf{Archetype} & \textbf{Provenance} & \textbf{Base score} &
\textbf{Envelope (\si{\degreeCelsius})} & \textbf{Confidence} &
\textbf{Entries} \\
\midrule
\multicolumn{6}{@{}l}{\textit{Green}}\\
Street tree canopy            & v1.1    & 75 & 0.5\,--\,3.0 & High   & 15 \\
Continuous shaded corridor    & v1.1    & 80 & 0.5\,--\,3.0 & Medium &  2 \\
Small green area with trees   & v1.1    & 65 & 0.3\,--\,1.5 & High   &  6 \\
Dense tree canopy             & v1.1    & 90 & 1.0\,--\,3.0 & High   & 10 \\
Established park              & v1.1    & 70 & 0.5\,--\,2.0 & High   & 10 \\
Non-canopy vegetation         & new     & 50 & 0.0\,--\,1.5 & Medium &  7 \\
Productive tree canopy        & derived & 60 & 0.5\,--\,2.0 & Low    &  3 \\
Productive ground planting    & derived & 45 & 0.0\,--\,1.5 & Low    &  6 \\
\addlinespace
\multicolumn{6}{@{}l}{\textit{Building}}\\
Extensive green roof          & v1.1    & 45 & 0.1\,--\,1.0 & Medium &  7 \\
Green fa\c{c}ade / living wall & v1.1   & 50 & 0.3\,--\,2.0 & Low    & 11 \\
\addlinespace
\multicolumn{6}{@{}l}{\textit{Blue-green}}\\
Bioretention                  & v1.1    & 55 & 0.1\,--\,0.8 & Low    &  8 \\
Blue-green corridor           & v1.1    & 85 & 1.0\,--\,3.0 & Medium &  5 \\
Riparian restoration          & v1.1    & 85 & 1.0\,--\,3.0 & Medium &  6 \\
Large open water body         & new     & 75 & 0.1\,--\,3.0 & Medium &  3 \\
Small constructed water       & new     & 40 & 0.0\,--\,1.0 & Medium &  6 \\
\addlinespace
\multicolumn{6}{@{}l}{\textit{Hybrid}}\\
Permeable shaded hardscape    & v1.1    & 70 & 0.5\,--\,2.5 & Medium &  2 \\
Enclosed courtyard greening   & v1.1    & 60 & 0.3\,--\,1.5 & Low    &  1 \\
Vegetated shade structure     & new     & 55 & 0.0\,--\,2.5 & Low    &  2 \\
\midrule
\multicolumn{5}{@{}l}{\textbf{Total entries}} & \textbf{110} \\
\bottomrule
\end{tabular}
\end{table}

\begin{keynote}
\noindent\textbf{Twelve of the eighteen archetypes are the previous version's
typologies, unchanged} --- same envelopes, same base scores, same evidence
ratings, same citations. They are the only classes in the library with
individually retrieved literature behind them, and this release keeps that
evidence rather than trading it for catalogue names.

\textbf{No previously published cooling envelope moved in this release}, with
one exception: the newly retrieved large water body ships at a floor of
\qty{0.1}{\degreeCelsius} rather than the \qty{0.5}{\degreeCelsius} proposed at
design time, because source verification showed the field evidence did not
support the higher floor. \Cref{sec:water-body} tells that story in full,
because it is a story about how the methodology is meant to work rather than an
erratum to be buried.
\end{keynote}

\textbf{How the base cooling score of a new archetype was placed.} The base
score is a relative \numrange{0}{100} strength index, not a physical quantity,
and the twelve inherited scores are not a function of the envelope midpoint ---
bioretention sits at 55 with a \qty{0.45}{\degreeCelsius} midpoint while the
green roof sits at 45 with a \qty{0.55}{\degreeCelsius} midpoint. There is
therefore no formula to apply, and inventing one would have moved twelve
calibrated values in order to make six new ones look derived. Each new score is
instead placed against a \emph{named comparator already in the library}, with
the argument stated where the archetype is derived below and in
\Cref{app:typologies}.

\subsection{Two retirements}
\label{sec:retirements}

Two typologies published at earlier versions no longer name a selectable
intervention. Both are methodology changes rather than renames: each carried
cited values and each appeared in a golden scenario.

\paragraph{Schoolyard greening is retired.} The source document is explicit
that a schoolyard is a \emph{land-use context}, not an intervention. Keeping it
would have meant asking the user ``what is the land use here?'' and then
offering, as an intervention, the land use they had just described. A school
site is instead offered the real interventions that suit it --- \num{72} of the
\num{110} entries name a school among their land uses --- through the
availability matrix of \Cref{sec:availability}. The evidence that supported
the retired typology is the same evidence supporting the small green area with
trees, which the catalogue's green playground entry inherits; nothing was
discarded except the duplicated question.

\paragraph{The mixed NbS package is retired} in favour of real packages
(\Cref{sec:packages}). The finding it encoded --- that no retrieved source
quantifies super-additive cooling --- has not changed and is not re-argued.
What changed is that the tool no longer needs an opaque averaged card to
express it: a user now selects the actual components, each is scored and shown
individually, and the combination rules do the work the averaged card was
standing in for.

\subsection{How envelopes were derived}
\label{sec:derivation}

\Cref{fig:cooling-evidence} shows the principal cross-typology anchor: the
per-measure daytime central estimates reported by \textcite{keravec2026}, with
two contrasting estimates overlaid.

\begin{figure}[htbp]
\centering
% Per-measure cooling evidence comparison.
% Bars: daytime pedestrian-level AIR temperature reduction central estimates
%   from the umbrella review (keravec2026).
% Overlaid markers, where an alternative estimate exists for the same measure:
%   diamond  = nighttime central estimate from the same review;
%   circle   = empirical-only park meta-analytic mean (bowler2010).
% Every value on this axis is an air temperature. No surface-temperature or
% thermal-comfort value appears here.
\begin{tikzpicture}
\begin{axis}[
  font=\footnotesize,
  width=0.80\textwidth,
  height=92mm,
  xbar,
  bar width=8pt,
  xmin=0, xmax=3.5,
  xtick={0,0.5,...,3.5},
  xticklabel style={/pgf/number format/fixed, /pgf/number format/precision=1},
  xlabel={Daytime pedestrian-level air temperature reduction (\si{\degreeCelsius})},
  xlabel style={font=\footnotesize},
  ytick=data,
  symbolic y coords={
    {Green walls},
    {Permeable pavements},
    {Water bodies},
    {Urban forests},
    {Green areas},
    {Shading devices},
    {Street trees},
    {Parks},
    {Green roofs}},
  y dir=reverse,
  yticklabel style={font=\footnotesize},
  enlarge y limits=0.13,
  axis lines=left,
  axis line style={draw=rulegrey, line width=0.5pt},
  tick style={draw=rulegrey},
  xmajorgrids=true,
  grid style={draw=rulegrey!45, line width=0.3pt, dash pattern=on 1pt off 1.5pt},
  legend style={
    at={(0.5,-0.20)}, anchor=north, font=\scriptsize, legend columns=1,
    draw=rulegrey, fill=white, legend cell align=left, row sep=1.5pt},
  clip=false,
]

% --- daytime central estimates (air temperature) -------------------------
\addplot[
  xbar, draw=criterragreen, fill=criterragreen!75,
  nodes near coords,
  nodes near coords style={font=\scriptsize\color{criterraink},
    /pgf/number format/fixed, /pgf/number format/precision=1,
    /pgf/number format/zerofill, xshift=1.5pt},
] coordinates {
  (3.0,{Green walls})
  (2.9,{Permeable pavements})
  (2.1,{Water bodies})
  (2.1,{Urban forests})
  (2.0,{Green areas})
  (2.0,{Shading devices})
  (1.5,{Street trees})
  (1.3,{Parks})
  (1.1,{Green roofs})
};
\addlegendentry{Umbrella-review central estimate, daytime air temperature}

% --- nighttime central estimates, where the review reports one -----------
\addplot[
  only marks, mark=diamond*, mark size=2.6pt,
  mark options={draw=criterraink, fill=criterraink}, forget plot,
] coordinates {
  (1.4,{Green walls})
  (1.2,{Parks})
};
\addlegendimage{only marks, mark=diamond*, mark size=2.6pt,
  mark options={draw=criterraink, fill=criterraink}}
\addlegendentry{Same review, nighttime central estimate (reported for two measures)}

% --- empirical-only park meta-analytic mean ------------------------------
\addplot[
  only marks, mark=*, mark size=2.4pt,
  mark options={draw=criterraink, fill=white, line width=0.8pt}, forget plot,
] coordinates {(0.94,{Parks})};
\addlegendimage{only marks, mark=*, mark size=2.4pt,
  mark options={draw=criterraink, fill=white, line width=0.8pt}}
\addlegendentry{Empirical-only park meta-analytic mean, \SI{0.94}{\degreeCelsius}}

\end{axis}
\end{tikzpicture}

\caption[Per-measure cooling evidence]{Central estimates of daytime
pedestrian-level \emph{air} temperature reduction by measure, from the umbrella
review of 64 systematic reviews by \textcite{keravec2026}. Overlaid markers
show, where an alternative estimate exists for the same measure, the nighttime
central estimate from the same review and the empirical-only park meta-analytic
mean of \qty{0.94}{\degreeCelsius} reported by \textcite{bowler2010}. The gap
between the park bar (\qty{1.3}{\degreeCelsius}, modelling-inclusive) and the
open marker (\qty{0.94}{\degreeCelsius}, empirical only) is the modelling bias
discussed in \Cref{sec:modelling-bias}; the diamonds show how much smaller
nighttime effects are (\Cref{sec:limit-daytime}).}
\label{fig:cooling-evidence}
\end{figure}

Three principles governed the derivation of an envelope from these estimates.

\textbf{An envelope is not a prediction.} The lower bound represents a poorly
performing instance of the archetype; the upper bound represents a well-executed
instance under favourable conditions. Site adjustment moves an assessment
\emph{within} this envelope (\Cref{sec:cooling}); it does not extend it.

\textbf{Central estimates are ceilings, not midpoints.} Because published
central estimates already describe interventions that were built, studied, and
in many cases modelled under favourable assumptions, adopted upper bounds sit
at or below the synthesis central estimate rather than above it.

\textbf{Where sources disagree, the conservative value is adopted and the
disagreement is documented} rather than averaged away. Three archetypes ---
the green roof, the green fa\c{c}ade, and the large open water body --- rest on
genuinely conflicting evidence, and all three are treated below.

\subsection{Calibration decisions that warrant reviewer attention}
\label{sec:calibration}

Four decisions materially lowered the tool's headline numbers relative to a
naive reading of the literature. Each is stated here because each is
contestable.

\subsubsection{Modelling bias is corrected downward}
\label{sec:modelling-bias}

\textcite{keravec2026} synthesises 64 systematic reviews, many of them
dominated by simulation studies, and simulation systematically reports larger
cooling effects than field measurement. The magnitude of this bias is visible
within the tool's own source set. For parks, the modelling-inclusive synthesis
gives a central estimate of \qty{1.3}{\degreeCelsius}, while
\textcite{bowler2010}, whose meta-analysis is restricted to empirical studies,
gives \qty{0.94}{\degreeCelsius} --- roughly \SI{72}{\percent} of the
modelling-inclusive figure, for the one measure where both are available.

Adopted upper bounds therefore sit at or below synthesis central estimates
rather than above them. The clearest case is the street tree canopy: the
synthesis reports \qty{4.0(16)}{\degreeCelsius} for street trees in tropical
climates, but the adopted ceiling is \qty{3.0}{\degreeCelsius}. The tropical
figure carries a \qty{\pm 1.6}{\degreeCelsius} interval and describes
favourable tropical conditions; adopting it as a general ceiling would overstate
typical performance.

This correction is a judgment, not a calculation. A single ratio observed for a
single measure is not a bias correction factor, and the methodology does not
present it as one --- it is used to justify the \emph{direction} of the
adjustment, not its size. A reviewer who considers the correction too
aggressive, or not aggressive enough, is disputing exactly the right thing.

\subsubsection{Packages do not claim super-additivity}
\label{sec:super-additivity}

An earlier draft of the methodology proposed a \qty{3.5}{\degreeCelsius}
ceiling for a combined intervention package, on the intuition that combining
measures should outperform any single measure. That value was rejected, and the
rejection survives the move to real packages unchanged.

No retrieved source quantifies super-additive cooling from combining
interventions. \textcite{gunawardena2017} describes green and blue features as
offering synergistic benefits, but qualitatively and without quantification.
\textcite{keravec2026} reports no synthesis estimate for combined packages at
all; its best-performing individual measures cluster at
\SIrange{2.0}{3.0}{\degreeCelsius}.

A package's reported temperature is therefore \textbf{capped at its
best-evidenced component and never summed} (\Cref{sec:packages}). Claiming that
a package outperforms its strongest component would be an unsourced assumption
placed at exactly the point where a user is most likely to select it. The
package's genuine advantage is expressed where evidence supports it --- breadth
of co-benefits, which are unioned across components --- and not as additional
degrees.

\subsubsection{Green roofs and green fa\c{c}ades are the weakest inherited
points, and are labelled as such}
\label{sec:weak-points}

\paragraph{Green roof.} Three sources point in different directions.
\textcite{keravec2026} places green roofs at \qty{1.1}{\degreeCelsius}, the
lowest-performing green measure in the synthesis. \textcite{zolch2016} finds
green roof effects at pedestrian level \emph{negligible}.
\textcite{santamouris2014} reports a \SIrange{0.3}{3}{\kelvin} reduction in
average urban ambient temperature --- but the abstract states this explicitly
for green roofs \emph{applied on a city scale} in \emph{simulation studies}.

That last qualification decides the matter. A site-level assessment evaluates
one roof on one building, not mass deployment across a city, and the physics
differ: a green roof's thermal benefit is concentrated at roof level and inside
the building beneath, and it does not propagate to the street canyon in the way
city-scale albedo and evapotranspiration changes would. The
\SIrange{0.3}{3}{\kelvin} figure is therefore \textbf{not} used as a site-level
value. The adopted envelope of
\SIrange{0.1}{1.0}{\degreeCelsius} sits near the bottom of the published
spread, and the tool states in its output that green roof benefits are
principally building-level rather than street-level.

\paragraph{Green fa\c{c}ade.} Here the sources conflict outright.
\textcite{keravec2026} ranks green walls at the top of its air-temperature
table with a \qty{3.0}{\degreeCelsius} central estimate --- but with an
approximately fivefold spread across climate subzones
(\Cref{fig:climate-dependence}). \textcite{zolch2016} ranks green fa\c{c}ades
\emph{below} trees on thermal comfort, at \SIrange{5}{10}{\percent} PET
reduction against \SI{13}{\percent} for trees.

The two are not directly comparable, since one reports air temperature and the
other a comfort index; but the ranking reversal is real and needs an
explanation. The most plausible one is measurement position: near-wall
measurements capture a strong local effect that does not extend across a street
canyon. This is a hypothesis, not a finding from either source, and it is
labelled as such.

Given an unresolved conflict between sources and a fivefold climate spread, the
archetype receives a conservative envelope of
\SIrange{0.3}{2.0}{\degreeCelsius}, an evidence confidence rating of
\emph{low}, and an explicit caveat in the tool's output. Its confidence rating
also caps the cooling-block confidence at medium regardless of how complete the
user's inputs are (\Cref{sec:evidence-propagation}). Eleven catalogue entries
inherit it, which makes it the second most widely inherited class in the
library and the one whose weakness reaches the most results.

\paragraph{And one further weak point.} Bioretention deserves the same flag. No
retrieved source quantifies air-temperature cooling for rain gardens or
bioswales specifically. \textcite{keravec2026} offers no bioswale-specific
estimate, and its neighbouring categories --- green areas at
\qty{2.0}{\degreeCelsius}, water bodies at \qty{2.1}{\degreeCelsius} --- are
not transferable to a small drainage feature.
\textcite{gunawardena2017} identifies tree-dominated greenspace as delivering
heat-stress relief and does not identify small blue features as significant
coolers. The adopted envelope of \SIrange{0.1}{0.8}{\degreeCelsius} is a
reasoned lower bound from adjacent evidence, flagged low confidence, and the
tool reports this archetype's primary benefits as stormwater management and
biodiversity rather than cooling.

\subsubsection{Newer, broader evidence is deliberately not adopted wholesale}
\label{sec:non-adoption}

\textcite{kumar2024} is a systematic review screening \num{27486} papers to
\num{202} analysed, covering \num{51} green--blue infrastructure types. It
post-dates the calibration of the twelve inherited archetypes, which makes it
the sharpest available external test of them, and it plays two distinct roles
in this methodology.

\textbf{As a validation}, it reports street trees at up to
\qty{2.8}{\degreeCelsius} by in-situ monitoring --- inside the tool's existing
\SIrange{0.5}{3.0}{\degreeCelsius} envelope. Newer and broader evidence
therefore supports the conservative calibration rather than overturning it,
which is the outcome a conservative calibration is supposed to produce and is
worth recording when it happens.

\textbf{As a non-adoption}, the same review reports means well above the tool's
envelopes for other types: green walls at
\qty{4.1(42)}{\degreeCelsius}, vegetated balconies at
\qty{3.8(27)}{\degreeCelsius}, and botanical gardens at
\qty{5.0(35)}{\degreeCelsius}. These are \textbf{not} adopted. Taking them
selectively would break internal consistency --- vegetated balconies would
outscore street trees --- and adopting them wholesale is a full recalibration of
the library, which would require its own sensitivity analysis and its own
version. It is recorded here and in the decision log as a candidate for a
future methodology review, not smuggled into this one. Note also that the
green-wall figure's standard deviation exceeds its mean.

A reviewer who believes the tool should adopt these figures is disputing
something specific and answerable, which is why the non-adoption is stated
rather than left implicit in a set of numbers that simply happen to be lower.

\subsection{The four archetypes retrieved for this version}
\label{sec:new-archetypes}

Each of the four exists because the catalogue contains entries that would
otherwise have inherited an envelope belonging to a different body of study. In
two cases that inheritance would have overstated the entry roughly threefold.

\subsubsection{The large open water body, and one envelope that moved at
implementation}
\label{sec:water-body}

The large open water body archetype --- inherited by urban lake restoration,
living shorelines, and coastal ecological corridors --- ships at
\SIrange{0.1}{3.0}{\degreeCelsius}, not the
\SIrange{0.5}{3.0}{\degreeCelsius} proposed when this release was designed.

\begin{keynote}
\noindent\textbf{A design may be approved; a citation is only settled once it
has been read.}

The approved design attributed a cooling magnitude near
\qty{2}{\degreeCelsius} to \textcite{yao2023a}. Source verification, carried out
before the value was allowed into configuration, established that the paper
reports no such figure: it measures a \SIrange{0.1}{0.6}{\degreeCelsius}
daytime cool island at urban lakes in two humid subtropical cities --- an order
of magnitude below the attributed value --- alongside a
\SIrange{1.2}{1.3}{\degreeCelsius} \emph{heat} island at night.

Holding the floor at \qty{0.5}{\degreeCelsius} would have made the tool claim
more than its most direct field evidence supports, which is precisely the
inflation the clipping rule of \Cref{sec:cooling} exists to prevent, and it
would have shipped a citation that did not say what the configuration claimed it
said. The floor was therefore lowered to the measured minimum.

Four further citation corrections were made in the same verification pass, none
of which moved a value: an author list, a nocturnal-warming figure conflated
between two papers by the same first author, a study of two parks under drought
that had been described as a study of allotment gardens, and a
mean-radiant-temperature range quoted more roundly than the source states. The
verification policy of \Cref{app:verification} is what caught all five, and it
is reported here rather than silently repaired because a methodology that only
publishes its successful checks is not publishing a check at all.
\end{keynote}

The resulting envelope spans a factor of thirty, and the width is the finding
rather than a defect. A remote-sensing-inclusive meta-analysis of 27 studies
gives a pooled \qty{2.5}{\kelvin} (\SIrange{1.9}{3.2}{\kelvin} at
\SI{95}{\percent}) \parencite{volker2013}, while direct field measurement of
actual urban lakes gives \SIrange{0.1}{0.6}{\degreeCelsius}
\parencite{yao2023a}; and \textcite{ampatzidis2020} identifies remote sensing
as the overestimating measurement mode. A narrower envelope would have to pick
a side that the evidence does not support picking, so the tool reports the
disagreement and its output caveat says so in as many words rather than leaving
a reader to wonder why the range is so wide.

The ceiling is unchanged at \qty{3.0}{\degreeCelsius}, on
\textcite{keravec2026}'s \qty{2.1}{\degreeCelsius} for water bodies and the
upper bound of the \textcite{volker2013} confidence interval. The base score of
75 is placed below the blue-green corridor at 85, because a water body alone
lacks the corridor's canopy component --- the corridor's envelope rests on the
canopy evidence as much as on the water evidence --- and because its
field-measured daytime performance is weak.

\subsubsection{Small constructed water, and the interpolation the engine refuses}
\label{sec:small-water}

Constructed wetlands, retention ponds, detention basins, water squares,
floating wetlands, and roof wetlands would naturally have inherited the
blue-green corridor at \SIrange{1.0}{3.0}{\degreeCelsius} on family resemblance
--- they are water, and water cools. Three independent lines say otherwise for
\emph{small constructed} water.

\textcite{jacobs2020} simulates 16 representative urban water bodies and reports
afternoon air temperature in surrounding spaces reduced by typically
\qty{0.2}{\degreeCelsius}, maximum \qty{0.6}{\degreeCelsius}; synthesising 20
further papers, 14 of them (\SI{70}{\percent}) report an effect of
\qty{1}{\degreeCelsius} or less, and the paper concludes that the local thermal
effect of small water bodies ``can be considered negligible in design
practice''. \textcite{yao2023b} measures an urban pond at
\qty{0.6}{\degreeCelsius} of daytime cooling against a reference site, with
\qty{1.8}{\degreeCelsius} of nocturnal \emph{warming}.
\textcite{ampatzidis2020} adds that sensible cooling can be offset by increased
water-vapour content. The adopted envelope is
\SIrange{0.0}{1.0}{\degreeCelsius}, with a genuine zero floor on
\textcite{jacobs2020}'s own conclusion, and a base score of 40 --- below the
extensive green roof at 45, because the two share a ceiling near
\qty{1.0}{\degreeCelsius} but this class has a true zero floor and the roof does
not. Without the split, six catalogue entries would have been overstated
roughly threefold.

\begin{keynote}
\noindent\textbf{The two water archetypes are never interpolated between by
area.} They are separate bodies of study, not two points on a published
dose--response curve: the size-threshold literature for water bodies is
land-surface-temperature only, and no published curve connects the two on air
temperature.

An area-based interpolation would look eminently reasonable and would be
unsourced. The engine therefore does not perform one, and the test suite asserts
its absence by scoring the same constructed wetland on two sites differing in
area by a factor of \num{2500} and requiring an identical ceiling, so that a
future edit cannot introduce it quietly.
\end{keynote}

\subsubsection{Non-canopy vegetation, and the surface/air conflation}
\label{sec:non-canopy}

Planting beds, groundcover, pocket meadows, roof meadows, and green medians ---
seven entries --- inherit a class of their own rather than the tree-based
classes, on the strength of one study that measures both metrics on the same
plots. \textcite{armson2012} found that grass reduced maximum \emph{surface}
temperature by up to \qty{24}{\degreeCelsius} while having little effect upon
\emph{globe} temperature, whereas shading reduced globe temperature by
\SIrange{5}{7}{\degreeCelsius}, and concludes that grass ``has little effect
upon local air or globe temperatures \dots\ whereas tree shade can provide
effective local cooling''.

That contrast is decisive precisely because it is measured on the same plots: a
\qty{24}{\degreeCelsius} surface reduction alongside almost no globe-temperature
change is not a small effect measured badly, it is a large effect on the wrong
quantity. It is also the mechanism by which most impressive published figures
for grass cooling are generated, and the tool's output caveat tells the user
exactly that.

The floor is a genuine zero, on \textcite{gill2013}, which models grass cooling
as reduced or lost in summer droughts when soils dry out, and on
\textcite{kraemer2022}, whose field measurement under the 2018--2019 Central
European drought found the grass-dominated park to be the hottest area during
the daytime while the tree-dominated park stayed consistently cooler. The
adopted envelope is \SIrange{0.0}{1.5}{\degreeCelsius} with a base score of 50,
below the small green area with trees at 65, which shares the
\qty{1.5}{\degreeCelsius} ceiling but adds the canopy shade that
\textcite{armson2012} shows is doing the work.

\subsubsection{Vegetated shade structures, and a figure that does not transfer}
\label{sec:shade-structure}

Green pergolas and green waiting shelters take
\SIrange{0.0}{2.5}{\degreeCelsius} at \emph{low} confidence, from
\textcite{chafer2020} --- the only measured comparison of a vine-covered
pergola against an identical bare support, giving a maximum daytime difference
of \qty{2.5}{\degreeCelsius} for west-oriented pergolas in a Mediterranean
climate. That paper's abstract headlines ``up to \qty{5}{\degreeCelsius}'', but
that figure compares greenery against \emph{unshaded} conditions rather than
against the bare pergola; the like-for-like contrast is the one adopted, and the
discrepancy is recorded in \Cref{app:verification} so that a reviewer checking
the abstract does not read an error into the ceiling.

\begin{keynote}
\noindent\textbf{The ``shading devices \qty{2.0}{\degreeCelsius}'' figure from
\textcite{keravec2026} is deliberately not used here.} That category aggregates
area-scale shading, where the cooled air volume persists against advection; a
pergola shades roughly \qty{10}{\metre\squared}. Applying the figure would rate
a pergola above measured tree canopy, which \textcite{ouyang2024} puts at
\qty{1.42}{\degreeCelsius}.

The same figure \emph{is} used by the permeable shaded hardscape archetype,
where the scale matches. The distinction is scale, not scepticism about the
source, and the test suite asserts both halves of it.
\end{keynote}

\textcite{ouyang2024} also separates the two axes that matter for shade
structures, measuring an air temperature reduction of
\qty{1.42}{\degreeCelsius} for natural shade against
\qty{1.31}{\degreeCelsius} for artificial, alongside mean radiant temperature
reductions of \qty{15.93}{\degreeCelsius} and \qty{13.71}{\degreeCelsius}
respectively --- an order of magnitude apart. Users will feel substantially
cooler under a pergola than the air-temperature figure suggests, and the output
caveat says so: under-selling an intervention is as misleading as over-selling
it, only in the other direction. \textcite{colter2019} supplies the ordering
that keeps the base score at 55, well below the permeable shaded hardscape at
70: sparse and open canopies group with solid constructed shade rather than
with dense canopy.

\subsection{The two derived archetypes, and a citable absence}
\label{sec:derived-archetypes}

Retrieval established that no class-specific cooling evidence exists for
productive landscapes. Both productive archetypes are therefore \emph{derived}
--- bounded from an adjacent class with the bounding argument stated --- and
both are rated \emph{low} confidence for that reason.

What makes this different from the bioretention case, where no source was found
at all, is that here \textbf{the absence itself is citable}.
\textcite{kumar2024} places allotments and city farms among the least-studied
of its \num{51} infrastructure types, with insufficient data to quantify their
cooling. The methodology can therefore state the gap as a finding rather than
as a search failure, which is a materially stronger position.

\textbf{Productive tree canopy} --- urban orchards, food forests, agroforestry
systems --- takes \SIrange{0.5}{2.0}{\degreeCelsius} at a base score of 60,
bounded \emph{below} dense tree canopy on two grounds: orchard and agroforestry
canopies are deliberately kept open for access and yield, so they sit below the
\SI{40}{\percent} canopy-cover inflection identified by \textcite{ziter2019},
and \textcite{colter2019} establishes that sparse canopies do not perform like
closed ones. The envelope matches the established park at a lower evidence
rating: the same range, honestly labelled as inherited rather than measured.

\textbf{Productive ground-level planting} --- community gardens, allotments,
school gardens, urban farms, productive roofs, rooftop farms --- inherits the
non-canopy vegetation envelope of \SIrange{0.0}{1.5}{\degreeCelsius} at a base
score of 45, below non-canopy vegetation at 50, because cultivated ground-level
cover is periodically broken between crops.

\begin{keynote}
\noindent\textbf{The best allotment-garden measurement available cannot be used
for a value.} \textcite{rost2020} measured 13 Berlin allotment complexes at an
average \qty{2.7}{\kelvin} cooler than built-up areas --- \emph{at night}. The
restriction is in the study's title, abstract, and design, and this tool reports
daytime air temperature. The study is cited as evidence of what the
productive-landscape literature does and does not measure, and for no numeric
value.

\textbf{Irrigation is not credited as extra degrees.} \textcite{cheung2022}
documents irrigation as \emph{sustaining} performance through drought rather
than raising it, measuring agricultural cooling at
\SIrange{0.09}{0.43}{\degreeCelsius} against modelled potentials of
\SIrange{2.1}{2.5}{\degreeCelsius}, and predicting an increase in the daily
\emph{minimum} temperature as stored heat is released overnight. One nuance is
carried carefully: that measured-versus-modelled gap is specific to
\emph{agricultural field} studies, and the same paper measures
\SIrange{-2}{-1}{\degreeCelsius} in irrigated urban parks, much closer to the
modelled values. The methodology therefore does \emph{not} claim that models
overstate irrigation cooling in general.
\end{keynote}

\subsection{Suitability conditions, inherited with twenty-one overrides}
\label{sec:suitability}

Each archetype declares the conditions under which an intervention is
unsuitable: a minimum viable site area, a soil requirement, an irrigation
requirement, and any unsuitable climate zones. Catalogue entries
\textbf{inherit} these, with per-entry overrides only where the catalogue's own
description of the entry makes an inherited value plainly wrong.

Inventing a minimum area and a soil requirement for each of \num{110} entries
would have manufactured exactly the precision the archetype model exists to
avoid. There are \num{21} overrides in total, each listed with its reason in
\Cref{app:typologies}, and \textbf{no new number enters the methodology through
an override}: every replacement value --- \qtylist{20;50;200;500;1000;2000}
{\metre\squared} --- is one already present in the cited library, borrowed from
the archetype that established it. A test fails if the override list grows past
\num{25} entries, because a long list would mean the inheritance ruling had
quietly been abandoned.

The overrides fall into three groups, and the second is the one worth reading.
Area overrides correct a scale mismatch: the dense tree canopy's
\qty{2000}{\metre\squared} describes the woodland its \emph{cooling evidence}
comes from, not the footprint of a microforest on a constrained site. Soil
overrides in the roof group set the requirement to \emph{none}, because a roof's
growing medium is built rather than ground, and requiring ground soil on a roof
would flag every roof entry whose archetype happens to be a ground-level class.
And the four ground-rooted fa\c{c}ade entries override soil in the
\emph{stricter} direction, to \emph{limited}, because the catalogue defines them
as rooted in natural ground where the containerised living-wall archetype
requires nothing --- an override that removes an option rather than adding one.

\paragraph{Land-use context is not a second list.} The urban-context
sub-indicator of \Cref{tab:suitability-rules} compares the site's land use
against the entry's own land-use list, which is the availability matrix's list
(\Cref{sec:availability}). An entry is therefore ``in context'' exactly where it
is offered --- one list, one meaning, and no opportunity for two lists to drift
apart.

When a disqualifying condition is met, the tool computes the assessment
transparently but attaches a prominent \emph{``not suitable for this site''}
flag naming the reason, which is carried into the recommendation text and the
exported report. Silently returning a merely lower score for a physically
implausible intervention --- riparian restoration on a site with no water
feature, for instance --- would misrepresent the finding as a matter of degree
when it is a matter of kind.

Unsuitable entries nonetheless remain fully selectable, including entries the
availability matrix did not offer for this site. A user deliberately testing a
hypothesis should not be overridden by the tool, and the flag follows the result
through to the report regardless.

One suitability rule is worth singling out because it encodes a caveat that
would otherwise be lost in prose. \textcite{gunawardena2017} notes that poorly
designed bluespace may \emph{exacerbate} heat stress under oppressive humid
conditions. Rather than leaving this as a written warning, the methodology
encodes it as a climate-suitability penalty: blue-green archetypes are
downgraded to the \emph{moderate} condition level in tropical-wet zones
(\Cref{tab:climate-matrix}). A caveat that changes a number is a caveat that
gets read.
